%----------------------------------
\section{Executive Summary of a Position Paper}
\label{sec:Executive Summary}
%----------------------------------

\subsection{Purpose and Audience}

Carlson and Worrell's~\cite{carlson2026relational} \emph{Relational Futures} report presents
findings from an Indigenous-led mixed-methods study, combining an online survey of 36 respondents
with yarning circles involving 22 participants, into how Aboriginal and Torres Strait Islander
peoples encounter artificial intelligence in everyday and professional life. Its purpose is not
merely descriptive. The authors argue that mainstream AI discourse assumes technological uptake is
inevitable and treats risk management as the main task, whereas their findings show Aboriginal and
Torres Strait Islander people experience AI as a relational, cultural and political force, not a
neutral tool. The report therefore functions as an advocacy and policy document, addressed to
three overlapping audiences: government and regulators being urged to strengthen Indigenous
governance and regulate high-risk AI uses; industry, whose systems are built on data and
infrastructure they do not control; and researchers, who are positioned as needing Indigenous-led
rather than Indigenous-consulted methodologies.

\subsection{Challenges and Opportunities}

The paper's central challenge is that AI systems are frequently deployed in domains where
Aboriginal and Torres Strait Islander peoples are already structurally over-represented, without
the governance authority to shape how those systems use their data. One participant summarised
this as AI arriving with ``no accountability, no checks and balances, no responsibility.'' The
authors situate this concern explicitly against precedents such as Robodebt, where an automated
decision system caused documented harm at scale. A further challenge is environmental and
material: routine AI use draws on water and energy often sourced from Indigenous lands, extending
the harms beyond data extraction into resource extraction. The corresponding opportunity the
report identifies is Indigenous Data Sovereignty (IDS): governance arrangements in which Aboriginal
and Torres Strait Islander communities hold authority over how data about them is collected,
interpreted and used, rather than being consulted after systems are already designed.

\subsection{Case Study: Predictive Risk Modelling in Child Protection}

A concrete illustration of this dynamic, distinct from the AI companions, chatbots and generative
systems discussed in the paper itself, is predictive risk modelling (PRM) in Australian child
protection. PRM tools use historical administrative data --- prior contact with services,
notifications, substantiations --- to flag families as high risk of future child maltreatment. A
scoping review by Krakouer, Wu Tan and Parolini~\cite{krakouer2021analysing} found such models are
being developed and piloted in Australian child protection systems, yet no empirical research had
examined their specific implications for Indigenous Australians, despite Aboriginal and Torres
Strait Islander children being dramatically over-represented in the system already. This is not a
peripheral technical gap. Because the training data reflects who has historically been surveilled
and reported, rather than who has actually caused harm, the model risks reifying that surveillance
as objective risk~\cite{keddell2015ethics}: a family or community already subject to
disproportionate contact is scored as higher risk, triggering further contact, which then becomes
tomorrow's training data. The loop rewards past bias with future authority.

This mirrors exactly what Carlson and Worrell describe at the level of AI governance generally. The
decisions at stake --- potential removal of children --- are among the most consequential a state
can make, yet the Aboriginal and Torres Strait Islander families and communities affected have no
authority over the design of the system used to inform them. Krakouer et
al.~\cite{krakouer2021analysing} conclude that participatory model development, transparency and
Indigenous data sovereignty are prerequisites before such tools can be considered fair for use with
Indigenous Australians; this is the same solution architecture \emph{Relational Futures} calls for
across AI more broadly --- Indigenous leadership over design and deployment, not after-the-fact
consultation on a finished product.

The opportunity side is equally transferable. The Maiam nayri Wingara Indigenous Data Sovereignty
Collective and Australian Indigenous Governance Institute's~\cite{maiam2018communique}
Communiqu\'e asserts that communities retain the right to refuse participation in data processes
inconsistent with their own principles, a right with direct application to families' ability to
contest or refuse PRM-informed interventions rather than simply be scored by them. The CARE
Principles --- Collective Benefit, Authority to Control, Responsibility and Ethics --- offer a
practical governance mechanism for operationalising this in a child-protection PRM: requiring, for
example, that Aboriginal Community Controlled Organisations hold authority over which variables a
model can use and how its outputs are actioned~\cite{carroll2020care}. This is not a fringe
position: the Australian Human Rights Commission's~\cite{ahrc2021humanrights} \emph{Human Rights
and Technology Final Report} reached comparable conclusions at a national level, recommending
stronger protections and an AI Safety Commissioner precisely because high-risk domains such as
social security and policing-adjacent decision-making had been shown to cause harm without
adequate oversight.

\subsection{Conclusion}

Read against a case study in a different domain, \emph{Relational Futures}' central claim holds
up: AI governance failures are not primarily technical failures to be fixed with better accuracy
metrics. Wherever automated systems inform decisions in areas where Aboriginal and Torres Strait
Islander peoples already carry a disproportionate burden of state intervention, the more decisive
variable is who holds authority over the system, not how well it performs on a validation set.
